\section*{Chapter \ref{ch:b3:structural-analysis} --- Structural Analysis}

\begin{solution}{exo:b3:structural-analysis:1}
The \omterm{def:b3:structural-analysis:stereo}{pole} trends $40 + 180 = \ang{220}$ and plunges $90 - 30 = \ang{60}$. The
plane dips $90 - 20 = \ang{70}$ towards $315 - 180 = \ang{135}$: $135/70$.
\end{solution}

\begin{solution}{exo:b3:structural-analysis:2}
$\theta = \ang{60}$. Equal-angle: $10\tan\ang{30} = \qty{5.77}{cm}$;
equal-area: $10\sqrt2\sin\ang{30} = \qty{7.07}{cm}$.
\end{solution}

\begin{solution}{exo:b3:structural-analysis:3}
$\vec n_1 = (\sin\ang{30}, 0, \cos\ang{30})$, $\vec n_2 = (-\sin\ang{30}, 0,
\cos\ang{30})$, $\vec n_1\cdot\vec n_2 = -0.25 + 0.75 = 0.5$: the \omterm{def:b3:structural-analysis:stereo}{poles} are
\ang{60} apart, the interlimb angle is \ang{120}, an open fold.
\end{solution}

\begin{solution}{exo:b3:structural-analysis:4}
$\vec d = (\cos\ang{60}, 0, -\sin\ang{60}) = (0.5, 0, -0.866)$ and $\vec u =
0.707\,\vec s + 0.707\,\vec d = (0.354, 0.707, -0.612)$: trend
$\operatorname{atan2}(0.354, 0.707) = \ang{26.6}$, plunge $\arcsin0.612 =
\ang{37.8}$. The rake lies between \ang{-90} (normal) and 0 (sinistral): an
oblique normal fault with a sinistral component.
\end{solution}

\begin{solution}{exo:b3:structural-analysis:5}
\begin{gather*}
  \vec n_1 = (0.5567, 0.3214, 0.7660),\qquad \vec n_2 = (-0.4698, -0.1710, 0.8660),\\
  \vec n_1\times\vec n_2 = (0.4093, -0.8420, 0.0558),
\end{gather*}
of length $0.9379$. Its opposite, $-\vec n_1\times\vec n_2$, points down:
trend \ang{334.1}, plunge $\arcsin(0.0558/0.9379) = \ang{3.4}$. $\vec n_1\cdot\vec
n_2 = 0.3469$, \ang{69.7} between the normals, interlimb angle \ang{110.3}:
an open fold.
\end{solution}

\begin{solution}{exo:b3:structural-analysis:6}
Plunge above \ang{60} means $\theta < \ang{30}$. Equal-area: $r^2 =
2\sin^2\ang{15} = 0.134$, \qty{13.4}{\percent} of the net; equal-angle:
$\tan^2\ang{15} = 0.072$, \qty{7.2}{\percent}. On the sphere the cap holds
$1 - \cos\ang{30} = 0.134$ of the hemisphere: only the equal-area net shows
it with its true share, so \omterm{def:b3:structural-analysis:stereo}{poles} are counted on it.
\end{solution}

\begin{solution}{exo:b3:structural-analysis:7}
\begin{gather*}
  T_{xx} = \frac{2\sin^2\ang{20}}3 = 0.078,\qquad
  T_{zz} = \frac{2\cos^2\ang{20} + 1}3 = 0.922,\\
  T_{xz} = -\frac{\sin\ang{20}\cos\ang{20} - \sin\ang{20}\cos\ang{20}}3 = 0,
\end{gather*}
and every term with $y$ is 0. The eigenvalues are $0.922$ (vertical),
$0.078$ (east) and $0$ (north): the \omterm{def:b3:structural-analysis:pi}{beta-axis} is horizontal, north--south.
The three \omterm{def:b3:structural-analysis:stereo}{poles} lie in one plane, so none has a component along its normal,
and $\vec u^{\mathsf T}T\vec u = 0$ for $\vec u$ along it.
\end{solution}

\begin{solution}{exo:b3:structural-analysis:8}
$a = 2$, $b = 1.5$, $k = 1/0.5 = 2$: constriction. $Z = 1/(1.8\times1.1) =
0.505$, $a = 1.636$, $b = 2.178$, $k = 0.636/1.178 = 0.54$: flattening.
\end{solution}

\begin{solution}{exo:b3:structural-analysis:9}
$S_1/S_2 = 2.25$ and $S_1S_2 = 1$: $S_1 = 1.5$, $S_2 = 0.667$; the long axis
was extended by \qty{50}{\percent}. The void has the semi-axes $1.5\times0.4 =
\qty{0.60}{mm}$ and $0.667\times0.4 = \qty{0.27}{mm}$.
\end{solution}

\begin{solution}{exo:b3:structural-analysis:10}
Equal-angle: $1/(2\cos^2\ang{30}) = 0.667$ along the radius and
$\tan\ang{30}/\sin\ang{60} = 0.577/0.866 = 0.667$ around: equal, conformal.
Equal-area: $\mathrm dr/(R\,\mathrm d\theta) = (\sqrt2/2)\cos\ang{30} = 0.612$
and $\sqrt2\sin\ang{30}/\sin\ang{60} = 0.707/0.866 = 0.816$; the product is
$0.500$, the area factor of the theorem. A small circle becomes an oval of axes
in the ratio $0.816/0.612 = 1.33$, long around the centre.
\end{solution}

\begin{solution}{exo:b3:structural-analysis:11}
The fault is vertical, so its normal is horizontal and perpendicular to the
vertical $\sigma_2$; by part (b) of the theorem the slip is perpendicular to
$\sigma_2$, horizontal, for every $\Phi$. With $\sigma_1 = 1$ along north,
$\sigma_3 = 0$ along east: $S\vec n = (0, -0.342, 0)$, $\vec n\cdot S\vec n =
0.117$, $\vec\tau = (0, -0.342, 0) - 0.117\,(0.940, -0.342, 0) = (-0.110,
-0.302, 0)$ and $-\vec\tau = (0.110, 0.302, 0)$, towards \ang{020}. The block
on the side of $\vec n$, the east side, moves north-north-east along the
strike: seen from the west side, it moves to the left, a sinistral fault (rake
0). $\sigma_2$ did not enter the computation.
\end{solution}

\begin{solution}{exo:b3:structural-analysis:12}
$S = (7.5 - 6.0)/1.0 = \qty{1.5}{km}$; $h = 6.0/1.5 = \qty{4.0}{km}$ below the
lower horizon. The section was $20 + 1.5 = \qty{21.5}{km}$ long, and changed
by $-1.5/21.5 = \qty{-7.0}{\percent}$.
\end{solution}

\begin{solution}{pb:b3:structural-analysis:1}
\textbf{1.} The \omterm{def:b3:structural-analysis:stereo}{pole} trends \ang{255} and plunges \ang{77}; $\theta =
\ang{13}$ and $r = 10\sqrt2\sin\ang{6.5} = \qty{1.60}{cm}$ from the centre,
towards \ang{255}.

\textbf{2.} $\vec n = (\sin\ang{13}\sin\ang{75}, \sin\ang{13}\cos\ang{75},
\cos\ang{13}) = (0.2173, 0.0582, 0.9744)$. Six nearly parallel unit vectors
add up to nearly 6: $|\vec S_E| = 5.9974$.

\textbf{3.} East limb: $\bar\alpha = \operatorname{atan2}(1.2583, 0.2017) =
\ang{80.9}$ and $\bar\delta = \arccos(5.8604/5.9974) = \ang{12.3}$; $\bar R
= 5.9974/6 = 0.99956$ and the scatter is $\sqrt{2\times0.00044} =
\qty{0.0297}{rad}$, or \ang{1.7}. West limb: $\bar\alpha =
\operatorname{atan2}(-1.2445, 0.1886) = \ang{278.6}$ and $\bar\delta =
\arccos(5.8636/5.9972) = \ang{12.1}$; $\bar R = 0.99953$, scatter \ang{1.8}.

\textbf{4.}
\[
  \vec S_E\times\vec S_W = (0.0774, -14.6714, 0.4883),
\]
of length $14.680$. Its opposite points down: the trend is $\operatorname{atan2}(-0.0774, 14.6714) = \ang{359.7}$ and the
plunge $\arcsin(0.4883/14.680) = \ang{1.9}$.

\textbf{5.} The cosine of the angle between the mean normals is
$32.835/35.968 = 0.9129$, the dot product of the sums divided by their
lengths: \ang{24.1} between the normals, an interlimb angle of \ang{155.9}:
an open fold, close to gentle.

\textbf{6.}
\begin{gather*}
  \vec m_E = (0.2098, 0.0336, 0.9772),\qquad \vec m_W = (-0.2075, 0.0314, 0.9777),\\
  \vec m_E - \vec m_W = (0.4173, 0.0022, -0.0005):
\end{gather*}
a horizontal normal towards \ang{089.7}. The axial plane is vertical, striking
\ang{359.7}, practically north--south: the fold is upright.

\textbf{7.} The girdle is the plane whose \omterm{def:b3:structural-analysis:stereo}{pole} is the \omterm{def:b3:structural-analysis:pi}{beta-axis}: dip direction
\ang{179.7}, dip $90 - 1.9 = \ang{88.1}$, almost the east--west vertical
plane. Every limb contains the fold axis, so every normal is perpendicular to
it and lies in the plane perpendicular to it.

\textbf{8.} $0.04400 + 0.00144 + 0.95456 = 1.00000$. $\vec b^{\mathsf T}T\vec
b = 0.00038$, the smallest eigenvalue: $\vec b$ is, to the rounding, the
eigenvector $\vec e_3$, and the girdle through the two mean \omterm{def:b3:structural-analysis:stereo}{poles} is the
best-fit girdle of all twelve.

\textbf{9.} $\arcsin\sqrt{0.00038} = \arcsin0.0195 = \ang{1.1}$.

\textbf{10.} Along the girdle the \omterm{def:b3:structural-analysis:stereo}{poles} sit about \ang{12.2} from the vertical,
on either side; $\sin^2\ang{12.2} = 0.045$, the mean square of their
east--west components, which is $\lambda_2$.

\textbf{11.} Yes: the \omterm{def:b3:structural-analysis:stereo}{poles} stray from the girdle by \ang{1.1} on average, no
more than the \ang{1.7} scatter of the readings of one limb about their mean.

\textbf{12.} $\ang{1.1}/\sqrt{12} = \ang{0.3}$: the plunge is \ang{1.9} $\pm$
\ang{0.3}.

\textbf{13.} A rake of \ang{-90}: a pure normal fault, dip-slip. Anderson:
$\sigma_1$ vertical, $\sigma_3$ horizontal along the dip direction,
\ang{120}--\ang{300}, $\sigma_2$ horizontal along the strike,
\ang{030}--\ang{210}.

\textbf{14.} The normal of $120/60$ lies in the vertical plane trending
\ang{120}, which holds $\vec e_1$ and $\vec e_3$: $n_2 = 0$. By part (b) the
slip is perpendicular to $\sigma_2$, the strike, for every $\Phi$: it is
dip-slip. The fault tells nothing of $\Phi$; in the formula both $n_2$ and
$\vec e_2\cdot\vec w$, hence numerator and denominator, vanish.

\textbf{15.} $\vec n = (\sin\ang{55}\sin\ang{150}, \sin\ang{55}\cos\ang{150},
\cos\ang{55}) = (0.4096, -0.7094, 0.5736)$; $n_1 = 0.574$; $n_2 = 0.5\times
0.4096 + 0.866\times(-0.7094) = -0.410$.

\textbf{16.} $\vec e_1\cdot\vec w = 0.198$ and $\vec e_2\cdot\vec w = 0.5\times
0.771 + 0.866\times0.605 = 0.909$, so $\Phi = -(0.574\times0.198)/(-0.410\times
0.909) = 0.1137/0.3727 = 0.305$.

\textbf{17.} $(0.305 + 0.294 + 0.301 + 0.306)/4 = 0.30$. The rakes were
rounded to the degree; one degree of rake moves $\Phi$ by about 0.02 on these
faults, so the half-degree of rounding explains the spread.

\textbf{18.} $\sigma_2 - \sigma_3 = 0.3(\sigma_1 - \sigma_3)$: $\sigma_2$ is
closer to $\sigma_3$ than to $\sigma_1$. Both horizontal stresses are well
below the vertical one; the region was extended mostly towards \ang{120}, a
little along \ang{030}.

\textbf{19.} With $\vec n = (0.750, -0.433, 0.500)$ and the slip $\vec u =
(0.433, -0.250, -0.866)$, $(\vec n - \vec u)/\sqrt2$ plunges \ang{75} towards
\ang{300}: P; $(\vec n + \vec u)/\sqrt2$ plunges \ang{15} towards \ang{120}:
T. They lie in the same vertical plane as $\sigma_1$ and $\sigma_3$ but
\ang{15} away from them, because the fault makes \ang{30}, not \ang{45}, with
$\sigma_1$.

\textbf{20.} $\tan a = \tan\ang{12.3}\times\cos(\ang{90} - \ang{80.9}) =
0.2180\times0.9874 = 0.2153$: $a = \ang{12.1}$.

\textbf{21.} $L_0 = 8000/\cos\ang{12} = \qty{8178.7}{m}$, $S =
\qty{178.7}{m}$, a relative change of $\cos\ang{12} - 1 =
\qty{-2.2}{\percent}$.

\textbf{22.} $4000\tan\ang{12} = \qty{850.2}{m}$; $A = \tfrac12\times8000\times
850.2 = \qty{3.40e6}{m^2}$.

\textbf{23.} $h = \num{3.40e6}/178.7 = \qty{19.0}{km}$. A detachment at the
base of the cover would give $h$ of a few hundred metres and an excess area
some fifty times smaller: the folds are not detached there. Either the
basement folded with the cover, or the synclines sank below the original
level, which makes $A$ smaller; the section alone cannot tell.

\textbf{24.} Heave $200/\tan\ang{60} = \qty{115.5}{m}$, slip $200/\sin\ang{60}
= \qty{230.9}{m}$.

\textbf{25.} $\tan a = \tan\ang{60}\cos\ang{30} = 1.500$: $a = \ang{56.3}$, and
the apparent heave is $200/1.5 = \qty{133.3}{m}$. The slip of the fault is
down the dip, towards \ang{120}: the section along \ang{120} shows true
lengths and should be the one restored.

\textbf{26.} The fault cuts the folds, so it is younger and is restored first,
moving the hanging wall back up the fault by \qty{230.9}{m}; the folds are
then unfolded. The \omterm{def:b3:structural-analysis:restoration}{pin lines} go on the axial surfaces of the synclines, in the
footwall, where the symmetric folds had no slip between the beds.

\textbf{27.} The Varden fold axis plunges
\textbf{\ang[text-series-to-math=true]{1.9}} towards the north (trend
\ang{359.7}): the \ang{2} built into the region, recovered to the
\ang{0.3} that twelve readings allow.
\end{solution}
